\chapter{Embolectomy and Thrombectomy}

\section*{Overview}
Embolectomy and thrombectomy are procedures that remove a blood clot from a blocked artery or vein, restoring blood flow to the affected organ or limb.

\section*{Alternative Names}
Clot removal; mechanical thrombectomy; pulmonary embolectomy; surgical embolectomy

\section*{Principle}
A catheter-based device or surgical approach captures or fragments the clot. Removing the obstruction restores perfusion and prevents irreversible tissue damage.

\section*{History}
Surgical embolectomy was pioneered by Trendelenburg in 1908, and modern catheter-based thrombectomy devices have expanded treatment of stroke and venous and arterial occlusion.

\section*{Why This Test or Procedure Is Ordered}
Performed for acute arterial occlusion of a limb or viscera, large-vessel occlusion stroke, and massive or selected submassive pulmonary embolism.

\section*{Is This a Primary or Secondary Test or Procedure}
Primary treatment for acute large-vessel occlusion.

\section*{How This Test or Procedure Is Performed}
Performed under imaging guidance (fluoroscopy or angiography). A catheter enters the vessel, and aspiration, stent-retrieval, or fragmentation removes the clot; open surgical embolectomy may be used for acute limb ischemia.

\section*{What Preparation Is Needed}
Imaging (CT angiography or angiography), assessment of symptom onset time, and correction of coagulopathy.

\section*{Contraindications and Precautions}
Hemorrhagic stroke (for stroke thrombectomy) and severe coagulopathy are contraindications.

\section*{Risks}
Reperfusion injury, bleeding, vessel injury or dissection, distal embolization, and contrast nephropathy.

\section*{What Happens After the Test or Procedure}
Perfusion is monitored; anticoagulation is often continued and the cause of embolism is investigated.

\section*{Understanding Your Results}
Restored flow is confirmed by angiography; clinical recovery depends on the duration of ischemia.

\section*{Normal Results}
Complete restoration of vessel patency and tissue perfusion.

\section*{Follow-up Tests and Procedures}
Imaging surveillance, anticoagulation, and treatment of the embolic source.

\section*{References}
\begin{enumerate}
  \item MedlinePlus. Embolectomy. U.S. National Library of Medicine; 2025.
  \item Current Medical Diagnosis and Treatment. McGraw-Hill; 2025.
\end{enumerate}

\clearpage
